pdfoutput=1
% Options for packages loaded elsewhere
\PassOptionsToPackage{unicode}{hyperref}
\PassOptionsToPackage{hyphens}{url}
\pdfoutput=1
\documentclass[
  12pt,
]{article}
\usepackage{xcolor}
\usepackage{amsmath,amssymb}
\setcounter{secnumdepth}{-\maxdimen} % remove section numbering
\usepackage{iftex}
\ifPDFTeX
  \usepackage[T1]{fontenc}
  \usepackage{textcomp} % provide euro and other symbols
\else % if luatex or xetex
  \usepackage{unicode-math} % this also loads fontspec
  \defaultfontfeatures{Scale=MatchLowercase}
  \defaultfontfeatures[\rmfamily]{Ligatures=TeX,Scale=1}
\fi
\usepackage{lmodern}
\ifPDFTeX\else
  % xetex/luatex font selection
\fi
% Use upquote if available, for straight quotes in verbatim environments
\IfFileExists{upquote.sty}{\usepackage{upquote}}{}
\IfFileExists{microtype.sty}{% use microtype if available
  \usepackage[]{microtype}
  \UseMicrotypeSet[protrusion]{basicmath} % disable protrusion for tt fonts
}{}
\makeatletter
\@ifundefined{KOMAClassName}{% if non-KOMA class
  \IfFileExists{parskip.sty}{%
    \usepackage{parskip}
  }{% else
    \setlength{\parindent}{0pt}
    \setlength{\parskip}{6pt plus 2pt minus 1pt}}
}{% if KOMA class
  \KOMAoptions{parskip=half}}
\makeatother
\usepackage{longtable,booktabs,array}
\usepackage{calc} % for calculating minipage widths
% Correct order of tables after \paragraph or \subparagraph
\usepackage{etoolbox}
\makeatletter
\patchcmd\longtable{\par}{\if@noskipsec\mbox{}\fi\par}{}{}
\makeatother
% Allow footnotes in longtable head/foot
\IfFileExists{footnotehyper.sty}{\usepackage{footnotehyper}}{\usepackage{footnote}}
\makesavenoteenv{longtable}
\setlength{\emergencystretch}{3em} % prevent overfull lines
\providecommand{\tightlist}{%
  \setlength{\itemsep}{0pt}\setlength{\parskip}{0pt}}
\usepackage{bookmark}
\IfFileExists{xurl.sty}{\usepackage{xurl}}{} % add URL line breaks if available
\urlstyle{same}
\hypersetup{
  hidelinks,
  pdfcreator={LaTeX via pandoc}}

\author{SCX}
\date{2026}

\begin{document}

\section{SCX Theory: Full Proof Chain
Audit}\label{scx-theory-full-proof-chain-audit}

\begin{quote}
\textbf{Date}: 2026-06-28 \textbar{} \textbf{Status}: Pre-submission
audit \textbf{Purpose}: Cross-check every theorem dependency, verify no
circular dependencies, identify hidden assumptions. \textbf{Scope}:
Theorems 1-3, Spring-1 (SE-1), Spring-2 (SE-2), and all connecting
propositions.
\end{quote}

\begin{center}\rule{0.5\linewidth}{0.5pt}\end{center}

\subsection{1. Dependency Graph
(Verified)}\label{dependency-graph-verified}

\begin{verbatim}
                         ┌──────────────────────┐
                         │   Theorem 3           │
                         │  (Unidentifiability)  │
                         │                       │
                         │   "Without A1-A6,     │
                         │    noise ≅ difficulty"│
                         └──────────┬───────────┘
                                    │
                    JUSTIFIES (why A1-A6 are needed)
                    No logical dependency on Thm 1
                                    │
                                    v
               ┌─────────────────────────────────────┐
               │   Assumption Set A1-A6               │
               │   (Disjoint training, independence,  │
               │    bounded loss, uniform noise,      │
               │    state homogeneity, balanced errors)│
               └─────────────────────────────────────┘
                    │                     │
                    v                     v
   ┌────────────────────────┐  ┌────────────────────────┐
   │   Theorem 1            │  │   Theorem 2            │
   │  (Noise Detection)     │  │  (Weak Feature)        │
   │                        │  │                        │
   │   "With A1-A6, F1→1    │  │   "Even with A1-A6,   │
   │    exponentially in M" │  │    if I(φ;S) ≤ δ,     │
   │                        │  │    SCX ≤ baseline      │
   │   PROOF: Hoeffding +   │  │    + O(√δ)"           │
   │   Chernoff on e_m      │  │                        │
   │                        │  │   PROOF: Pinsker/BH +  │
   └───────────┬────────────┘  │   Fano + data process. │
               │               └───────────┬────────────┘
               │                           │
               │  provides S_0             │  provides boundary
               │  (consensus scores)       │  O(√δ) on advantage
               │                           │
               v                           v
   ┌────────────────────────────────────────────────────┐
   │   Spring-1: Theorem SE-1                           │
   │  (Convergence of Self-Evolution)                   │
   │                                                    │
   │   "(S_t, θ_t) → (S*, θ*) under C1-C9"              │
   │                                                    │
   │   PROOF: Lyapunov descent (Theorem 12.5) +         │
   │   memory bank stabilization (Lemma SE-1.2) +       │
   │   vanishing displacement (Lemma SE-1.3) +          │
   │   fixed-point characterization (Lemma SE-1.4)      │
   │                                                    │
   │   DEPENDS ON:                                      │
   │   - Thm 1: provides noise detection signal for S_0 │
   │   - Thm 2: characterizes when advantage vanishes   │
   │   - C1'-C9: structural conditions                  │
   │   - Theorem 12.5: Lyapunov descent (proven)        │
   └───────────────────────┬────────────────────────────┘
                           │
                           │  provides Lyapunov structure
                           │  and convergence guarantee
                           v
   ┌────────────────────────────────────────────────────┐
   │   Spring-2: Theorem SE-2                           │
   │  (Completeness Bound)                              │
   │                                                    │
   │   "Finite-time termination under physical limits"  │
   │                                                    │
   │   PROOF: Finite configuration space (Prop SE-3) +  │
   │   Lyapunov descent → no cycles of length > 1 +     │
   │   bounded improvement → finite steps to fixed point│
   │                                                    │
   │   DEPENDS ON:                                      │
   │   - Spring-1: Lyapunov descent + convergence       │
   │   - Physical constraints: finite data, finite      │
   │     precision, finite parameterization             │
   └────────────────────────────────────────────────────┘
\end{verbatim}

\textbf{Verdict}: The dependency graph is a \textbf{DAG} (directed
acyclic graph). \textbf{No circular dependencies found.}

\begin{center}\rule{0.5\linewidth}{0.5pt}\end{center}

\subsection{2. Dependency-by-Dependency
Verification}\label{dependency-by-dependency-verification}

\subsubsection{2.1 Chain: Thm 3 → A1-A6 → Thm
1}\label{chain-thm-3-a1-a6-thm-1}

\textbf{Claim}: Theorem 1's assumptions (A1-A6) are justified by Theorem
3's unidentifiability result.

\textbf{Direction}: Thm 3 → Thm 1 (justification, not logical
dependency)

\textbf{Verification}:

\begin{longtable}[]{@{}
  >{\raggedright\arraybackslash}p{(\linewidth - 4\tabcolsep) * \real{0.2800}}
  >{\raggedright\arraybackslash}p{(\linewidth - 4\tabcolsep) * \real{0.3200}}
  >{\raggedright\arraybackslash}p{(\linewidth - 4\tabcolsep) * \real{0.4000}}@{}}
\toprule\noalign{}
\begin{minipage}[b]{\linewidth}\raggedright
Check
\end{minipage} & \begin{minipage}[b]{\linewidth}\raggedright
Result
\end{minipage} & \begin{minipage}[b]{\linewidth}\raggedright
Evidence
\end{minipage} \\
\midrule\noalign{}
\endhead
\bottomrule\noalign{}
\endlastfoot
Does Thm 1's proof require Thm 3? & \textbf{NO} --- Thm 1 is standalone
& Thm 1 proof uses only A1-A6 + Hoeffding/Chernoff; Thm 3 is never cited
in the proof \\
Does Thm 3's proof require Thm 1? & \textbf{NO} --- Thm 3 is standalone
& Thm 3 proof is a constructive counterexample; references Thm 1 only
for A1-A6 definitions \\
Is the relationship correctly characterized? & \textbf{YES} & Document
00: ``Theorem 3\ldots establishes why assumptions are required. Theorem
1\ldots quantifies how well SCX works when its assumptions hold.'' \\
Are A1-A6 consistent between the two theorems? & \textbf{YES} & Same
A1-A6 used in both; Thm 3's Section 4 maps each assumption to what it
breaks \\
\end{longtable}

\textbf{Verdict}: \textbf{PASS.} No circularity. The relationship is:
Thm 3 justifies \textbf{why} A1-A6 are needed; Thm 1 proves
\textbf{what} A1-A6 achieve. This is a conceptual dependency
(motivation), not a logical dependency (proof structure).

\textbf{Note on arrow direction}: Some documents describe this as ``Thm
1 → Thm 3'' (reading left-to-right as ``Thm 1's assumptions require Thm
3's justification''). The actual logical flow is Thm 3 → A1-A6 → Thm 1.
This is a harmless notational difference but should be clarified in the
arXiv submission.

\begin{center}\rule{0.5\linewidth}{0.5pt}\end{center}

\subsubsection{2.2 Chain: Thm 1 → Spring-1 (Noise Detection → Gatekeeper
Initialization)}\label{chain-thm-1-spring-1-noise-detection-gatekeeper-initialization}

\textbf{Claim}: Theorem 1's noise detection guarantee provides the
signal that initializes the gatekeeper \(S_0\) and continuously
calibrates \(S_t\).

\textbf{Direction}: Thm 1 → Spring-1 (input dependency)

\textbf{Verification}:

\begin{longtable}[]{@{}
  >{\raggedright\arraybackslash}p{(\linewidth - 4\tabcolsep) * \real{0.2800}}
  >{\raggedright\arraybackslash}p{(\linewidth - 4\tabcolsep) * \real{0.3200}}
  >{\raggedright\arraybackslash}p{(\linewidth - 4\tabcolsep) * \real{0.4000}}@{}}
\toprule\noalign{}
\begin{minipage}[b]{\linewidth}\raggedright
Check
\end{minipage} & \begin{minipage}[b]{\linewidth}\raggedright
Result
\end{minipage} & \begin{minipage}[b]{\linewidth}\raggedright
Evidence
\end{minipage} \\
\midrule\noalign{}
\endhead
\bottomrule\noalign{}
\endlastfoot
Does Spring-1 use Thm 1's consensus score? & \textbf{YES} & \(S_0\) is
initialized from \(C(x)\) (Thm 1's consensus score); \(S_t\) updates use
SCXUpdate which depends on consensus \\
Is Proposition SE-1.4 correctly stated? & \textbf{YES} & ``Throughout
the self-evolution process\ldots the noise detection guarantee of
Theorem 1 applies at each time step \(t\)'' \\
Is the F1 global aggregation correct? & \textbf{YES (FIXED)} &
DEFECT-01/02 fixed; Lemma F now uses correct global F1 via linear
FPR/FNR decomposition \\
Does Spring-1 add assumptions beyond Thm 1? & \textbf{YES} & C1-C9 are
additional (Lipschitz, RM rates, two-timescale, exploration, etc.) ---
but they are \textbf{orthogonal} to A1-A6, not contradictory \\
Could Spring-1 work without Thm 1? & \textbf{NO} & Without Thm 1's noise
detection signal, \(S_0\) has no meaningful initialization; the
gatekeeper would have no basis for acceptance/rejection \\
\end{longtable}

\textbf{Verdict}: \textbf{PASS.} Dependency is real and correctly
documented. Thm 1 provides the essential noise detection signal.
Spring-1 adds dynamical assumptions (C1-C9) that are consistent with but
independent of A1-A6.

\begin{center}\rule{0.5\linewidth}{0.5pt}\end{center}

\subsubsection{2.3 Chain: Thm 2 → Spring-1 (Weak Feature → Degradation
Rate)}\label{chain-thm-2-spring-1-weak-feature-degradation-rate}

\textbf{Claim}: Theorem 2's weak feature bound characterizes when SCX
cannot outperform the loss baseline, providing the boundary condition
for Spring-1's effectiveness.

\textbf{Direction}: Thm 2 → Spring-1 (boundary condition)

\textbf{Verification}:

\begin{longtable}[]{@{}
  >{\raggedright\arraybackslash}p{(\linewidth - 4\tabcolsep) * \real{0.2800}}
  >{\raggedright\arraybackslash}p{(\linewidth - 4\tabcolsep) * \real{0.3200}}
  >{\raggedright\arraybackslash}p{(\linewidth - 4\tabcolsep) * \real{0.4000}}@{}}
\toprule\noalign{}
\begin{minipage}[b]{\linewidth}\raggedright
Check
\end{minipage} & \begin{minipage}[b]{\linewidth}\raggedright
Result
\end{minipage} & \begin{minipage}[b]{\linewidth}\raggedright
Evidence
\end{minipage} \\
\midrule\noalign{}
\endhead
\bottomrule\noalign{}
\endlastfoot
Does Spring-1 depend on Thm 2 for correctness? & \textbf{NO} ---
Spring-1's convergence proof does not use Thm 2 & Spring-1 (Theorem
SE-1) proof in Document 06 uses only C1-C9 + Lyapunov descent; Thm 2 is
not cited in the convergence proof itself \\
Does Thm 2 constrain Spring-1's effectiveness? & \textbf{YES} & When
\(I(\phi; S) \to 0\), SCX cannot outperform baseline; this means
Spring-1 converges to a fixed point that may be no better than the
baseline \\
Is this relationship correctly documented? & \textbf{YES} & Document 06
Section 14: ``Thm 2\ldots bounds the quality of the fixed point: if
features are weak, even the optimal fixed point cannot outperform the
loss baseline'' \\
Is the \(O(\sqrt{\delta})\) bound consistent with Spring-1's rate? &
\textbf{PARTIALLY} & Thm 2 gives \(O(\sqrt{\delta})\) bound on SCX
advantage. Spring-1 gives \(O(\alpha_t + \beta_t)\) convergence rate.
These are orthogonal rates (one is feature-quality bound, one is
optimization rate). No contradiction. \\
\end{longtable}

\textbf{Verdict}: \textbf{PASS.} Thm 2 provides a \textbf{boundary
condition} for Spring-1 --- it tells us when convergence is meaningful
--- but does not enter Spring-1's proof logically. This is correct: a
boundary theorem should not be required for the convergence proof
itself.

\begin{center}\rule{0.5\linewidth}{0.5pt}\end{center}

\subsubsection{2.4 Chain: Spring-1 → Spring-2 (Convergence →
Completeness)}\label{chain-spring-1-spring-2-convergence-completeness}

\textbf{Claim}: Theorem SE-1's convergence guarantee provides the
Lyapunov structure that Theorem SE-2 uses to prove finite-time
termination.

\textbf{Direction}: Spring-1 → Spring-2 (logical dependency)

\textbf{Verification}:

\begin{longtable}[]{@{}
  >{\raggedright\arraybackslash}p{(\linewidth - 4\tabcolsep) * \real{0.2800}}
  >{\raggedright\arraybackslash}p{(\linewidth - 4\tabcolsep) * \real{0.3200}}
  >{\raggedright\arraybackslash}p{(\linewidth - 4\tabcolsep) * \real{0.4000}}@{}}
\toprule\noalign{}
\begin{minipage}[b]{\linewidth}\raggedright
Check
\end{minipage} & \begin{minipage}[b]{\linewidth}\raggedright
Result
\end{minipage} & \begin{minipage}[b]{\linewidth}\raggedright
Evidence
\end{minipage} \\
\midrule\noalign{}
\endhead
\bottomrule\noalign{}
\endlastfoot
Does SE-2's proof require SE-1's convergence? & \textbf{YES} & SE-2
(Document 07) proof: ``By the strict Lyapunov property (Theorem SE-1,
assumption), \(\Phi(q_{t+1}) < \Phi(q_t)\) unless
\(q_{t+1} = q_t\).'' \\
Does SE-2 add assumptions beyond SE-1? & \textbf{YES} & SE-2 adds
physical constraints: finite data (\(N_{\max} < \infty\)), finite
precision (\(\varepsilon_{\text{mach}} > 0\)), finite parameterization.
These are \textbf{independent} of SE-1's C1-C9. \\
Is SE-2's finite-configuration argument valid? & \textbf{YES (FIXED)} &
DEFECT-15 replaced ``finite \(\mathcal{X}\)'' with covering-number
argument; Proposition SE-3 now correctly bounds
\(|\mathcal{Q}| < \infty\) under physical constraints \\
Does SE-2's fixed-point characterization match SE-1's? & \textbf{YES} &
Both characterize
\(S^* = \text{SCXUpdate}(S^*, M_\infty, f_{\theta^*})\) and
\(\nabla L_{S^*}(\theta^*) = 0\) \\
Could SE-2 hold without SE-1? & \textbf{PARTIALLY} & SE-2's
finite-configuration argument (termination in finite time) holds
regardless of SE-1, but the \textbf{quality} of the fixed point (that
it's a Lyapunov minimum) depends on SE-1 \\
\end{longtable}

\textbf{Verdict}: \textbf{PASS.} Spring-2 logically depends on Spring-1
for the Lyapunov structure. Spring-2's additional physical constraints
are independent and verifiable.

\begin{center}\rule{0.5\linewidth}{0.5pt}\end{center}

\subsubsection{2.5 Cross-Verification: Hidden
Assumptions}\label{cross-verification-hidden-assumptions}

We check for assumptions that are used but not explicitly stated in the
dependency chain.

\begin{longtable}[]{@{}
  >{\raggedright\arraybackslash}p{(\linewidth - 4\tabcolsep) * \real{0.6047}}
  >{\raggedright\arraybackslash}p{(\linewidth - 4\tabcolsep) * \real{0.1860}}
  >{\raggedright\arraybackslash}p{(\linewidth - 4\tabcolsep) * \real{0.2093}}@{}}
\toprule\noalign{}
\begin{minipage}[b]{\linewidth}\raggedright
Hidden Assumption Check
\end{minipage} & \begin{minipage}[b]{\linewidth}\raggedright
Result
\end{minipage} & \begin{minipage}[b]{\linewidth}\raggedright
Details
\end{minipage} \\
\midrule\noalign{}
\endhead
\bottomrule\noalign{}
\endlastfoot
Does Thm 1 assume expert models are ``good enough''
(\(\mu_s < (K-1)/K\))? & \textbf{Explicit} & Lemma 1 states this as the
separation condition; Corollary 3 (uniform detectability) formalizes
it \\
Does Thm 2 assume the clustering algorithm (k-means) achieves the Fano
lower bound? & \textbf{FIXED (DEFECT-11)} & Now explicitly states it's
conservative; uses single-sample \(\delta\) which gives an upper bound
on SCX performance \\
Does Spring-1 assume the memory bank stabilizes? & \textbf{Explicit} &
Lemma SE-1.2 proves this under C1 (covering dimension) + monotonicity \\
Does Spring-1 assume the student loss is smooth? & \textbf{Explicit} &
Condition C2 (Lipschitz student) + \(L_g\)-smoothness derived from C2 \\
Does Spring-2 assume the system is deterministic? & \textbf{Explicit} &
``The self-evolution update is a deterministic mapping
\(G: \mathcal{Q} \to \mathcal{Q}\)'' \\
Does Thm 3's construction use A1-A2 (which it's trying to justify)? &
\textbf{NO (verified)} & Thm 3 Section 2.4 Step 4 explicitly states ``we
construct the experts to be conditionally independent\ldots this does
not assume A1-A2 hold universally'' --- it's part of the constructive
counterexample \\
\end{longtable}

\textbf{Verdict}: \textbf{No hidden assumptions found.} All assumptions
are explicitly stated in their respective theorem documents. One
borderline case (Thm 3 using conditional independence in its
construction) is correctly handled --- the construction demonstrates
that even when A1-A2 hold, the counterexample exists, which only
strengthens the theorem.

\begin{center}\rule{0.5\linewidth}{0.5pt}\end{center}

\subsection{3. Notation Consistency
Audit}\label{notation-consistency-audit}

\subsubsection{3.1 Critical Notation
Checks}\label{critical-notation-checks}

\begin{longtable}[]{@{}
  >{\raggedright\arraybackslash}p{(\linewidth - 12\tabcolsep) * \real{0.1290}}
  >{\raggedright\arraybackslash}p{(\linewidth - 12\tabcolsep) * \real{0.1129}}
  >{\raggedright\arraybackslash}p{(\linewidth - 12\tabcolsep) * \real{0.1129}}
  >{\raggedright\arraybackslash}p{(\linewidth - 12\tabcolsep) * \real{0.1129}}
  >{\raggedright\arraybackslash}p{(\linewidth - 12\tabcolsep) * \real{0.1613}}
  >{\raggedright\arraybackslash}p{(\linewidth - 12\tabcolsep) * \real{0.1613}}
  >{\raggedright\arraybackslash}p{(\linewidth - 12\tabcolsep) * \real{0.2097}}@{}}
\toprule\noalign{}
\begin{minipage}[b]{\linewidth}\raggedright
Symbol
\end{minipage} & \begin{minipage}[b]{\linewidth}\raggedright
Thm 1
\end{minipage} & \begin{minipage}[b]{\linewidth}\raggedright
Thm 2
\end{minipage} & \begin{minipage}[b]{\linewidth}\raggedright
Thm 3
\end{minipage} & \begin{minipage}[b]{\linewidth}\raggedright
Spring-1
\end{minipage} & \begin{minipage}[b]{\linewidth}\raggedright
Spring-2
\end{minipage} & \begin{minipage}[b]{\linewidth}\raggedright
Consistent?
\end{minipage} \\
\midrule\noalign{}
\endhead
\bottomrule\noalign{}
\endlastfoot
\(\mathcal{S}\) & State space & State space & State space &
\textbf{Gatekeeper scores \(S_t\)} & Gatekeeper scores &
\textbf{CONFLICT} --- documented in 00\_notation \\
\(S\) & State variable & State variable & State variable &
\textbf{Gatekeeper scoring function} & Gatekeeper function &
\textbf{CONFLICT} --- documented \\
\(K\) & \(|\mathcal{Y}|\) (classes) & N/A & \(|\mathcal{Y}|\) & N/A &
N/A & \textbf{CONSISTENT} \\
\(K_S\) & Not used & \(|\mathcal{S}|\) (states) & Not used &
\(|\mathcal{S}|\) & \(|\mathcal{S}|\) & \textbf{CONSISTENT} (but added
in polished versions) \\
\(\eta\) & Noise rate & Noise rate & Noise rate & N/A & N/A &
\textbf{CONSISTENT} \\
\(\delta\) & Confidence parameter & \textbf{Mutual information bound} &
N/A & N/A & N/A & \textbf{CONFLICT} (different meanings in Thm 1 vs Thm
2) \\
\(\Delta_s\) & Separation gap & Not used & Not used & Not used & Not
used & \textbf{UNIQUE to Thm 1} \\
\(\rho_s\) & State probability & Not used & State probability & Not used
& Not used & \textbf{CONSISTENT} \\
\(M\) & Expert count & Expert count & Expert count & Not used & Not used
& \textbf{CONSISTENT} \\
\(M_t\) & Not used & Not used & Not used & \textbf{Memory bank} & Memory
bank & \textbf{UNIQUE to Spring-1/2} \\
\(\Phi\) & Not used & Feature space & Not used & \textbf{Lyapunov
function} & Lyapunov function & \textbf{CONFLICT} (feature space vs
Lyapunov) \\
\(\phi\) & Not used & Feature map & Not used & Not used & Not used &
\textbf{UNIQUE to Thm 2} \\
\end{longtable}

\subsubsection{3.2 Notation Conflicts Requiring
Attention}\label{notation-conflicts-requiring-attention}

\textbf{Conflict 1: \(\mathcal{S}\) (state space) vs.~\(S_t\)
(gatekeeper scoring function)} - \textbf{Affected documents}: Theorems
1-3 use \(\mathcal{S}\) for state space. Spring-1/2 use \(S_t\) for
gatekeeper scores. - \textbf{Resolution}: Document
00\_notation\_and\_dependencies.md clarifies this. The polished theorem
documents (polished/) use \(K_S = |\mathcal{S}|\) for state space
cardinality. - \textbf{Risk}: Low. Context disambiguates (Thm 1-3 never
mention \(S_t\); Spring-1/2 always subscript \(S_t\) with time index). -
\textbf{Recommendation}: Add a one-sentence note in Spring-1:
``Notation: \(\mathcal{S}\) denotes the state space (as in Thm 1-3);
\(S_t\) denotes the gatekeeper scoring function at time \(t\).''

\textbf{Conflict 2: \(\delta\) (confidence in Thm 1 vs.~mutual
information in Thm 2)} - \textbf{Affected documents}: Thm 1 uses
\(\delta\) for confidence parameter (e.g., \(1-\delta\)). Thm 2 uses
\(\delta\) for \(I(\phi; S) \leq \delta\). - \textbf{Risk}: Medium. Both
are standard in their respective contexts but could confuse a reader who
reads both theorems sequentially. - \textbf{Recommendation}: Rename Thm
2's \(\delta\) to \(\delta_{\phi}\) or \(\iota\) (iota for information)
in the arXiv version.

\textbf{Conflict 3: \(\Phi\) (feature space in Thm 2 vs.~Lyapunov
function in Spring-1/2)} - \textbf{Affected documents}: Thm 2 uses
\(\Phi\) for feature space. Spring-1/2 use \(\Phi\) for the Lyapunov
function. - \textbf{Risk}: Medium. The contexts are very different (Thm
2: information theory; Spring: dynamical systems) but the symbol appears
in both. - \textbf{Recommendation}: Spring documents should use
\(\mathcal{V}\) or \(\Psi\) for the Lyapunov function instead of
\(\Phi\), or Thm 2 should use \(\mathcal{F}_{\phi}\) for the feature
space.

\subsubsection{3.3 Notation Conflicts:
Assessment}\label{notation-conflicts-assessment}

\begin{longtable}[]{@{}
  >{\raggedright\arraybackslash}p{(\linewidth - 4\tabcolsep) * \real{0.3704}}
  >{\raggedright\arraybackslash}p{(\linewidth - 4\tabcolsep) * \real{0.2593}}
  >{\raggedright\arraybackslash}p{(\linewidth - 4\tabcolsep) * \real{0.3704}}@{}}
\toprule\noalign{}
\begin{minipage}[b]{\linewidth}\raggedright
Severity
\end{minipage} & \begin{minipage}[b]{\linewidth}\raggedright
Count
\end{minipage} & \begin{minipage}[b]{\linewidth}\raggedright
Examples
\end{minipage} \\
\midrule\noalign{}
\endhead
\bottomrule\noalign{}
\endlastfoot
\textbf{Critical} (proof-breaking) & 0 & --- \\
\textbf{Moderate} (reader-confusing) & 3 & \(\mathcal{S}\) vs \(S_t\),
\(\delta\) (dual use), \(\Phi\) (dual use) \\
\textbf{Minor} (stylistic) & 2 & \(K\) vs \(K_S\) (polished versions use
\(K_S\)), \(\rho_s\) (Thm 1) vs \(\rho\) (Thm 3's state probability) \\
\end{longtable}

\textbf{Overall}: Notation is \textbf{mostly consistent} within each
theorem's scope. Cross-theorem conflicts are documented and resolvable
with minor renaming.

\begin{center}\rule{0.5\linewidth}{0.5pt}\end{center}

\subsection{4. Dependency Strength
Assessment}\label{dependency-strength-assessment}

\subsubsection{4.1 Logical Dependencies
(Proof-Critical)}\label{logical-dependencies-proof-critical}

\begin{longtable}[]{@{}
  >{\raggedright\arraybackslash}p{(\linewidth - 6\tabcolsep) * \real{0.1395}}
  >{\raggedright\arraybackslash}p{(\linewidth - 6\tabcolsep) * \real{0.0930}}
  >{\raggedright\arraybackslash}p{(\linewidth - 6\tabcolsep) * \real{0.2093}}
  >{\raggedright\arraybackslash}p{(\linewidth - 6\tabcolsep) * \real{0.5581}}@{}}
\toprule\noalign{}
\begin{minipage}[b]{\linewidth}\raggedright
From
\end{minipage} & \begin{minipage}[b]{\linewidth}\raggedright
To
\end{minipage} & \begin{minipage}[b]{\linewidth}\raggedright
Strength
\end{minipage} & \begin{minipage}[b]{\linewidth}\raggedright
If broken, what fails?
\end{minipage} \\
\midrule\noalign{}
\endhead
\bottomrule\noalign{}
\endlastfoot
C1-C9 & Spring-1 & \textbf{STRONG} & Convergence proof fails without
these conditions \\
Lyapunov descent (Thm 12.5) & Spring-1 & \textbf{STRONG} & Without
Lyapunov descent, only Cesàro-mean convergence remains (Theorem 10.1) \\
Spring-1 & Spring-2 & \textbf{STRONG} & SE-2's fixed-point quality
guarantee depends on SE-1's Lyapunov structure \\
A1-A6 & Thm 1 & \textbf{STRONG} & Noise detection guarantee fails
without these assumptions \\
Fano + Pinsker/BH & Thm 2 & \textbf{STRONG} & Weak feature bound depends
on these inequalities \\
\end{longtable}

\subsubsection{4.2 Conceptual Dependencies
(Motivation/Interpretation)}\label{conceptual-dependencies-motivationinterpretation}

\begin{longtable}[]{@{}
  >{\raggedright\arraybackslash}p{(\linewidth - 6\tabcolsep) * \real{0.2143}}
  >{\raggedright\arraybackslash}p{(\linewidth - 6\tabcolsep) * \real{0.1429}}
  >{\raggedright\arraybackslash}p{(\linewidth - 6\tabcolsep) * \real{0.3214}}
  >{\raggedright\arraybackslash}p{(\linewidth - 6\tabcolsep) * \real{0.3214}}@{}}
\toprule\noalign{}
\begin{minipage}[b]{\linewidth}\raggedright
From
\end{minipage} & \begin{minipage}[b]{\linewidth}\raggedright
To
\end{minipage} & \begin{minipage}[b]{\linewidth}\raggedright
Strength
\end{minipage} & \begin{minipage}[b]{\linewidth}\raggedright
Purpose
\end{minipage} \\
\midrule\noalign{}
\endhead
\bottomrule\noalign{}
\endlastfoot
Thm 3 & A1-A6 & \textbf{MODERATE} & Justifies why A1-A6 are not
arbitrary \\
Thm 1 & Spring-1 (\(S_0\)) & \textbf{MODERATE} & Provides the initial
gatekeeper signal \\
Thm 2 & Spring-1 (boundary) & \textbf{WEAK} & Characterizes when
convergence is meaningful \\
\end{longtable}

\subsubsection{4.3 Critical Path Analysis}\label{critical-path-analysis}

The \textbf{critical proof path} is:

\begin{verbatim}
A1-A6 → Thm 1 (noise detection signal)
C1-C9 → Theorem 12.5 (Lyapunov descent) → Spring-1 (convergence) → Spring-2 (completeness)
Thm 2 (boundary condition)
\end{verbatim}

If any link in the strong dependency chain breaks, the corresponding
theorem's proof is invalid. Currently:

\begin{itemize}
\tightlist
\item
  \textbf{All strong dependency links verified intact.}
\item
  \textbf{Thm 1}: A1-A6 are explicitly stated and used. Proof is
  complete (Hoeffding + Chernoff).
\item
  \textbf{Theorem 12.5}: Requires reference-set replay mechanism. Proof
  is complete (importance sampling + reference-based SCXUpdate).
\item
  \textbf{Spring-1}: Depends on Theorem 12.5 for Lyapunov descent. With
  Theorem 12.5, the proof is complete.
\item
  \textbf{Spring-2}: Depends on Spring-1 for Lyapunov structure. Proof
  is complete under physical constraints.
\end{itemize}

\begin{center}\rule{0.5\linewidth}{0.5pt}\end{center}

\subsection{5. Gap Inventory (What Is NOT
Proven)}\label{gap-inventory-what-is-not-proven}

\begin{longtable}[]{@{}
  >{\raggedright\arraybackslash}p{(\linewidth - 8\tabcolsep) * \real{0.1509}}
  >{\raggedright\arraybackslash}p{(\linewidth - 8\tabcolsep) * \real{0.2453}}
  >{\raggedright\arraybackslash}p{(\linewidth - 8\tabcolsep) * \real{0.1887}}
  >{\raggedright\arraybackslash}p{(\linewidth - 8\tabcolsep) * \real{0.1887}}
  >{\raggedright\arraybackslash}p{(\linewidth - 8\tabcolsep) * \real{0.2264}}@{}}
\toprule\noalign{}
\begin{minipage}[b]{\linewidth}\raggedright
Gap ID
\end{minipage} & \begin{minipage}[b]{\linewidth}\raggedright
Description
\end{minipage} & \begin{minipage}[b]{\linewidth}\raggedright
Location
\end{minipage} & \begin{minipage}[b]{\linewidth}\raggedright
Severity
\end{minipage} & \begin{minipage}[b]{\linewidth}\raggedright
Resolution
\end{minipage} \\
\midrule\noalign{}
\endhead
\bottomrule\noalign{}
\endlastfoot
G1 & Lyapunov descent WITHOUT reference-set replay & 10\_lyapunov, Thm
12.2 & \textbf{FUNDAMENTAL} (proven impossible) & Theorem 12.2 proves
impossibility; reference-set replay (Thm 12.5) closes the gap \\
G2 & Infinite \(\mathcal{X}\) extension of Lemma SE-1.2 &
06\_fixed\_point, Section 14 & \textbf{MINOR} & Covering-number argument
(DEFECT-15) resolves for \(\varepsilon\)-precision; full functional
convergence is open \\
G3 & Phase transitions between regimes 1-4 & 06\_fixed\_point, Section
11 & \textbf{CONJECTURED} & Not needed for main theorems; documented as
open problem \\
G4 & Perpetual discovery rate bound & 06\_fixed\_point, Proposition
SE-1.2 & \textbf{CONJECTURED} & Not needed for main convergence
result \\
G5 & Thm 3's \(K>2\) construction uses ``completely random experts'' &
03\_unidentifiability, Appendix A & \textbf{MINOR} (construction is
valid but extreme) & Acknowledged as extreme case; \(K=2\) construction
is the practically relevant one \\
\end{longtable}

\begin{center}\rule{0.5\linewidth}{0.5pt}\end{center}

\subsection{6. Final Verdict}\label{final-verdict}

\subsubsection{6.1 Audit Results}\label{audit-results}

\begin{longtable}[]{@{}
  >{\raggedright\arraybackslash}p{(\linewidth - 2\tabcolsep) * \real{0.5789}}
  >{\raggedright\arraybackslash}p{(\linewidth - 2\tabcolsep) * \real{0.4211}}@{}}
\toprule\noalign{}
\begin{minipage}[b]{\linewidth}\raggedright
Criterion
\end{minipage} & \begin{minipage}[b]{\linewidth}\raggedright
Result
\end{minipage} \\
\midrule\noalign{}
\endhead
\bottomrule\noalign{}
\endlastfoot
\textbf{Circular dependencies} & \textbf{NONE FOUND.} The dependency
graph is a strict DAG. \\
\textbf{Hidden assumptions} & \textbf{NONE FOUND.} All assumptions are
explicitly stated and justified. \\
\textbf{Notation conflicts} & \textbf{3 MODERATE conflicts} identified
(\(\mathcal{S}\)/\(S_t\), \(\delta\), \(\Phi\)). All documented. None
are proof-breaking. \\
\textbf{Broken proof chains} & \textbf{NONE.} All strong dependencies
are verified intact. \\
\textbf{Unproven claims marked as proven} & \textbf{1 RETRACTED}
(SE-1.5, \(1/\sqrt{N_t}\) rate --- DEFECT-05). Corrected version now
acknowledges non-resolution. \\
\textbf{Gaps acknowledged} & \textbf{5 gaps} documented in Section 5.
All are either resolved, minor, or explicitly marked as conjectured. \\
\end{longtable}

\subsubsection{6.2 Readiness for arXiv}\label{readiness-for-arxiv}

The proof chain is \textbf{structurally sound} for arXiv submission with
the following caveats:

\begin{enumerate}
\def\labelenumi{\arabic{enumi}.}
\item
  \textbf{Theorem 12.5 (reference-set replay)} should be the primary
  Lyapunov result. The impossibility result (Theorem 12.2) provides
  strong motivation for the reference-set replay mechanism.
\item
  \textbf{Notation conflicts} should be resolved before submission (see
  Section 3.2 recommendations).
\item
  \textbf{Gap G2} (infinite \(\mathcal{X}\) functional convergence)
  should be explicitly acknowledged as an open problem.
\item
  \textbf{Thm 3 Appendix A} (\(K>2\) construction) should include the
  caveat that the construction uses extreme ``completely random
  experts'' and that the \(K=2\) case is the practically relevant one.
  This is already documented.
\end{enumerate}

\subsubsection{6.3 Theorem Status Summary}\label{theorem-status-summary}

\begin{longtable}[]{@{}
  >{\raggedright\arraybackslash}p{(\linewidth - 6\tabcolsep) * \real{0.1667}}
  >{\raggedright\arraybackslash}p{(\linewidth - 6\tabcolsep) * \real{0.1481}}
  >{\raggedright\arraybackslash}p{(\linewidth - 6\tabcolsep) * \real{0.3704}}
  >{\raggedright\arraybackslash}p{(\linewidth - 6\tabcolsep) * \real{0.3148}}@{}}
\toprule\noalign{}
\begin{minipage}[b]{\linewidth}\raggedright
Theorem
\end{minipage} & \begin{minipage}[b]{\linewidth}\raggedright
Status
\end{minipage} & \begin{minipage}[b]{\linewidth}\raggedright
Proof Completeness
\end{minipage} & \begin{minipage}[b]{\linewidth}\raggedright
Key Conditions
\end{minipage} \\
\midrule\noalign{}
\endhead
\bottomrule\noalign{}
\endlastfoot
\textbf{Theorem 1} (Noise Detection) & \textbf{PROVEN} & 100\% &
A1-A6 \\
\textbf{Theorem 2} (Weak Feature) & \textbf{PROVEN} & 100\% &
\(\delta\)-weak \(\phi\), state balance \\
\textbf{Theorem 3} (Unidentifiability) & \textbf{PROVEN} & 100\% &
Constructive counterexample \\
\textbf{Proposition 3} (State-Conditioned Weighting) & \textbf{PROVEN} &
100\% & Gibbs inequality \\
\textbf{Proposition 4} (Compression Fidelity) & \textbf{PROVEN} & 100\%
& A1-A4 (compression-specific) \\
\textbf{Theorem 12.5} (Lyapunov Descent with Replay) & \textbf{PROVEN} &
100\% & C1'-C9 + reference-set replay \\
\textbf{Theorem 12.2} (Lyapunov Impossibility without Replay) &
\textbf{PROVEN} & 100\% & Proves necessity of replay \\
\textbf{Spring-1 (SE-1)} (Convergence) & \textbf{PROVEN} (with Thm 12.5)
& 95\% & C1-C9 + Thm 12.5 \\
\textbf{Spring-2 (SE-2)} (Completeness) & \textbf{PROVEN} & 100\% &
Physical constraints + SE-1 \\
\textbf{Regime 3} (Perpetual Discovery) & \textbf{CONJECTURED} & 0\% &
Open problem \\
\textbf{Phase Diagram} & \textbf{CONJECTURED} & 0\% & Open problem \\
\end{longtable}

\begin{center}\rule{0.5\linewidth}{0.5pt}\end{center}

\emph{End of PROOF\_CHAIN\_AUDIT.md}

\begin{center}\rule{0.5\linewidth}{0.5pt}\end{center}

\subsection{References
(Audit-Specific)}\label{references-audit-specific}

\begin{enumerate}
\def\labelenumi{\arabic{enumi}.}
\tightlist
\item
  All theorem documents in \texttt{/theory/theorems/} and
  \texttt{/theory/theorems/polished/}
\item
  All self-evolution documents in \texttt{/theory/self\_evolution/}
\item
  Notation document:
  \texttt{/theory/theorems/polished/00\_notation\_and\_dependencies.md}
\item
  Definitions:
  \texttt{/theory/definitions/01\_state\_conditioned\_risk.md}
\item
  Propositions:
  \texttt{/theory/propositions/03\_state\_conditioned\_weighting.md},
  \texttt{/theory/propositions/04\_compression\_fidelity.md}
\end{enumerate}

\end{document}
